\section{Related Work}
\label{sec:related}

\paragraph{Memory poisoning attacks on LLM agents.}
\citet{chen2024agentpoison} introduce \agentpoison{}, the first attack to
exploit the retrieval mechanism of LLM agent memory systems.  They formulate
trigger optimization as a constrained gradient problem over the embedding
space of a dense retriever (DPR~\citep{karpukhin2020dense}) and demonstrate
attacks on three production agent frameworks (LangChain, Voyager, ReAct),
achieving $\asrr \geq 0.80$ at below 0.1\% poison rate.
\citet{dong2025minja} extend the threat model to a \emph{query-only} setting:
the attacker has no direct memory access and relies on the agent's auto-storage
of adversarial reasoning chains.  Their progressive-shortening indication
protocol achieves ISR of 98.2\% across three experimental platforms.
\citet{injecmem2026} propose a retriever-agnostic anchor strategy that achieves
broad recall via topic-covering passages designed to match many query types,
enabling injection in a single interaction.
\citet{memorygraft2025} demonstrate that successful experiences stored in
agent memory can be \emph{grafted}: an attacker implants a malicious
``experience record'' that is later retrieved and replicated as a legitimate
agent action, exploiting the union retrieval over lexical and embedding similarity.
Concurrent work on tool-use manipulation~\citep{yang2024watch} shows that
similar poisoning strategies transfer to agent tool selection, broadening
the attack surface beyond memory retrieval.
\citet{li2026memorypoisoning} provide an empirical comparison of memory
poisoning attacks and defenses in EHR agents, confirming that attack
effectiveness is largely robust to corpus size---consistent with our
corpus-size ablation (Appendix~\ref{sec:ablation_appendix}).

\paragraph{RAG corpus poisoning.}
Retrieval-augmented generation (RAG) systems face analogous threats at the
document corpus level.  \citet{zou2024poisonedrag} demonstrate that
adversarial documents optimized via HotFlip~\citep{ebrahimi2018hotflip}
can force specific LLM outputs when retrieved, achieving $> 90\%$ ASR
at million-scale knowledge bases.  \citet{chaudhari2024phantom}
construct phantom documents that match specific queries while appearing
semantically benign to human inspection.  \citet{xue2024badrag} show
that a small number of adversarially crafted documents can redirect LLM
answers across entire topic domains, consistent with our corpus-size ablation.
\citet{arzanipour2025ragsecurity} provide the first formal threat model
for RAG systems, defining adversary access tiers and formally bounding
the attack surface at the document and embedding levels.
\citet{tan2025revprag} detect poisoned RAG responses via LLM activation
analysis (ResNet18 on multi-layer activations), achieving $98\%$ TPR at
$1\%$ FPR---a complementary detection signal to our query-similarity approach.
Our work differs from RAG poisoning in that the adversarial entries are
\emph{agent-generated, semantically coherent memories} rather than external
documents, the threat persists across sessions without further adversarial
interaction, and defenses must operate at write time without full corpus visibility.

\paragraph{Backdoor attacks on neural retrieval.}
Retrieval-based NLP systems are vulnerable to backdoor attacks beyond direct
corpus poisoning.  \citet{yang2021rethinking} and \citet{carlini2024poisoning}
demonstrate that embedding-space backdoors can cause targeted retrievals
in open-domain question answering systems.  \citet{badchain2024} show that
chain-of-thought reasoning prompts can be backdoored, causing models to
follow malicious reasoning chains when a trigger phrase appears---analogous
to \minja{}'s bridging-step mechanism.
The ASB benchmark \citep{zhang2024asb} provides a unified evaluation
framework for backdoor attacks and defenses in classification settings;
we adapt its matrix format to the memory-agent threat model.
Neural Cleanse \citep{wang2019neuralcleanse} and STRIP \citep{gao2019strip}
are representative detection methods that inspired our proactive and composite
defense designs.  However, these methods are fundamentally model-centric:
Neural Cleanse requires white-box model access and reverse-engineers trigger
patterns in parameter space; STRIP perturbs \emph{inference-time inputs} and
detects Trojan activation via entropy collapse.  Neither applies to the
memory-poisoning setting: our defenses target raw \emph{stored content}
rather than model weights or inference activations, and must operate at
write time without access to the retrieval model's internal states.
This architectural difference motivates SAD's calibration-based approach,
which detects anomalous \emph{similarity} rather than anomalous \emph{activation}.

\paragraph{LLM watermarking for content provenance.}
\citet{zhao2024watermark} propose a provably robust token-level watermark
for LLM-generated text based on a green/red vocabulary partition, achieving
AUROC $> 0.97$ and providing formal guarantees under $\ell_\infty$-bounded
paraphrasing.  \citet{kirchenbauer2023watermark} independently develop
a similar scheme (KGW) and analyze its robustness to paraphrasing attacks.
\citet{hu2023unbiased} propose an unbiased watermarking scheme that avoids
the quality degradation associated with hard vocabulary partitions.
We adapt the Zhao et al.\ unigram scheme to character-level memory entries
that do not pass through a generative model (Section~\ref{sec:defenses});
our evasion evaluation mirrors the paraphrasing and dilution attacks studied
in \citet{kirchenbauer2023watermark} and extends them to copy-paste dilution
and adaptive vocabulary substitution.

\paragraph{Prompt injection and indirect injection attacks.}
\citet{perez2022ignore} first characterize prompt injection attacks, where
a crafted user input overrides system instructions.  \citet{greshake2023not}
extend this to \emph{indirect prompt injection}: adversarial content in
retrieved documents (web pages, emails, tool outputs) hijacks the agent's
behavior without the user's knowledge.  Memory poisoning is orthogonal:
the adversarial content is stored silently and triggers on future queries
without further adversarial interaction, persisting across sessions.
This makes memory poisoning strictly more dangerous in multi-session
deployments: a single poisoning event affects \emph{all} future interactions.

\paragraph{LLM agent security and memory systems.}
The LLM agent security landscape has broadened rapidly.
\citet{memgpt2023} introduce MemGPT's hierarchical memory system, which
our framework targets via the Mem0~\citep{mem0} and A-MEM~\citep{amem}
wrappers.  \citet{touvron2023llama} and \citet{openai2023gpt4} describe
production-scale LLMs that underpin modern agent deployments.
\citet{yang2024watch} identify that agent tool selection and memory retrieval
share a common embedding-based vulnerability; our attack-defense matrix
provides the first systematic cross-attack coverage in this setting.
\citet{gu2019badnets} establish the foundational backdoor threat model
that underlies our definition of memory poisoning (Definition~1).
\citet{zhang2024asb} benchmark 10 prompt injection attacks, memory poisoning,
and 11 defenses across 13 LLM backbones in the ASB framework; we adopt
its matrix format and extend coverage to pre-ingestion filtering defenses.
\citet{amemguard2025} propose a proactive defense combining consensus-based
validation over multiple memory retrievals with a ``lessons learned'' store;
their 95\% ASR reduction is complementary to our write-time filtering approach.

\paragraph{Anomaly detection for adversarial inputs.}
Our \sad{} defense draws on a line of anomaly detection work in NLP.
\citet{gao2019strip} use input perturbation to detect Trojan-activated
predictions; we adapt the key insight---that adversarial inputs cluster
differently from benign ones in representation space---to the memory entry
setting.  \citet{goldblum2022dataset} and \citet{hayase2021spectre} develop
spectral methods for detecting poisoned samples in supervised learning;
their methods require a labeled set of both clean and poisoned training examples
to compute a covariance structure and identify poisoned directions.
In the memory-agent setting, no labeled poison examples are available at
detection time, making spectral methods inapplicable.
\sad{} sidesteps this by exploiting the \emph{query-similarity signal}:
poison entries must be semantically close to victim queries (by attack
design), while benign entries are not, enabling calibration from benign
content alone.  \citet{lee2018simple} show that Mahalanobis distance to
class centroids is an effective out-of-distribution detector;
our cosine similarity threshold is the retrieval-space analogue, adapted
to the unsupervised, query-conditioned memory setting.
